% Options for packages loaded elsewhere
\PassOptionsToPackage{unicode}{hyperref}
\PassOptionsToPackage{hyphens}{url}
\documentclass[
]{article}
\usepackage{xcolor}
\usepackage{amsmath,amssymb}
\setcounter{secnumdepth}{-\maxdimen} % remove section numbering
\usepackage{iftex}
\ifPDFTeX
  \usepackage[T1]{fontenc}
  \usepackage[utf8]{inputenc}
  \usepackage{textcomp} % provide euro and other symbols
\else % if luatex or xetex
  \usepackage{unicode-math} % this also loads fontspec
  \defaultfontfeatures{Scale=MatchLowercase}
  \defaultfontfeatures[\rmfamily]{Ligatures=TeX,Scale=1}
\fi
\usepackage{lmodern}
\ifPDFTeX\else
  % xetex/luatex font selection
\fi
% Use upquote if available, for straight quotes in verbatim environments
\IfFileExists{upquote.sty}{\usepackage{upquote}}{}
\IfFileExists{microtype.sty}{% use microtype if available
  \usepackage[]{microtype}
  \UseMicrotypeSet[protrusion]{basicmath} % disable protrusion for tt fonts
}{}
\makeatletter
\@ifundefined{KOMAClassName}{% if non-KOMA class
  \IfFileExists{parskip.sty}{%
    \usepackage{parskip}
  }{% else
    \setlength{\parindent}{0pt}
    \setlength{\parskip}{6pt plus 2pt minus 1pt}}
}{% if KOMA class
  \KOMAoptions{parskip=half}}
\makeatother
\usepackage{color}
\usepackage{fancyvrb}
\newcommand{\VerbBar}{|}
\newcommand{\VERB}{\Verb[commandchars=\\\{\}]}
\DefineVerbatimEnvironment{Highlighting}{Verbatim}{commandchars=\\\{\}}
% Add ',fontsize=\small' for more characters per line
\newenvironment{Shaded}{}{}
\newcommand{\AlertTok}[1]{\textcolor[rgb]{1.00,0.00,0.00}{\textbf{#1}}}
\newcommand{\AnnotationTok}[1]{\textcolor[rgb]{0.38,0.63,0.69}{\textbf{\textit{#1}}}}
\newcommand{\AttributeTok}[1]{\textcolor[rgb]{0.49,0.56,0.16}{#1}}
\newcommand{\BaseNTok}[1]{\textcolor[rgb]{0.25,0.63,0.44}{#1}}
\newcommand{\BuiltInTok}[1]{\textcolor[rgb]{0.00,0.50,0.00}{#1}}
\newcommand{\CharTok}[1]{\textcolor[rgb]{0.25,0.44,0.63}{#1}}
\newcommand{\CommentTok}[1]{\textcolor[rgb]{0.38,0.63,0.69}{\textit{#1}}}
\newcommand{\CommentVarTok}[1]{\textcolor[rgb]{0.38,0.63,0.69}{\textbf{\textit{#1}}}}
\newcommand{\ConstantTok}[1]{\textcolor[rgb]{0.53,0.00,0.00}{#1}}
\newcommand{\ControlFlowTok}[1]{\textcolor[rgb]{0.00,0.44,0.13}{\textbf{#1}}}
\newcommand{\DataTypeTok}[1]{\textcolor[rgb]{0.56,0.13,0.00}{#1}}
\newcommand{\DecValTok}[1]{\textcolor[rgb]{0.25,0.63,0.44}{#1}}
\newcommand{\DocumentationTok}[1]{\textcolor[rgb]{0.73,0.13,0.13}{\textit{#1}}}
\newcommand{\ErrorTok}[1]{\textcolor[rgb]{1.00,0.00,0.00}{\textbf{#1}}}
\newcommand{\ExtensionTok}[1]{#1}
\newcommand{\FloatTok}[1]{\textcolor[rgb]{0.25,0.63,0.44}{#1}}
\newcommand{\FunctionTok}[1]{\textcolor[rgb]{0.02,0.16,0.49}{#1}}
\newcommand{\ImportTok}[1]{\textcolor[rgb]{0.00,0.50,0.00}{\textbf{#1}}}
\newcommand{\InformationTok}[1]{\textcolor[rgb]{0.38,0.63,0.69}{\textbf{\textit{#1}}}}
\newcommand{\KeywordTok}[1]{\textcolor[rgb]{0.00,0.44,0.13}{\textbf{#1}}}
\newcommand{\NormalTok}[1]{#1}
\newcommand{\OperatorTok}[1]{\textcolor[rgb]{0.40,0.40,0.40}{#1}}
\newcommand{\OtherTok}[1]{\textcolor[rgb]{0.00,0.44,0.13}{#1}}
\newcommand{\PreprocessorTok}[1]{\textcolor[rgb]{0.74,0.48,0.00}{#1}}
\newcommand{\RegionMarkerTok}[1]{#1}
\newcommand{\SpecialCharTok}[1]{\textcolor[rgb]{0.25,0.44,0.63}{#1}}
\newcommand{\SpecialStringTok}[1]{\textcolor[rgb]{0.73,0.40,0.53}{#1}}
\newcommand{\StringTok}[1]{\textcolor[rgb]{0.25,0.44,0.63}{#1}}
\newcommand{\VariableTok}[1]{\textcolor[rgb]{0.10,0.09,0.49}{#1}}
\newcommand{\VerbatimStringTok}[1]{\textcolor[rgb]{0.25,0.44,0.63}{#1}}
\newcommand{\WarningTok}[1]{\textcolor[rgb]{0.38,0.63,0.69}{\textbf{\textit{#1}}}}
\usepackage{longtable,booktabs,array}
\usepackage{caption}
\captionsetup[table]{skip=6pt}
\newcounter{none} % for unnumbered tables
\usepackage{calc} % for calculating minipage widths
% Correct order of tables after \paragraph or \subparagraph
\usepackage{etoolbox}
\makeatletter
\patchcmd\longtable{\par}{\if@noskipsec\mbox{}\fi\par}{}{}
\makeatother
\setlength{\emergencystretch}{3em} % prevent overfull lines
\providecommand{\tightlist}{%
  \setlength{\itemsep}{0pt}\setlength{\parskip}{0pt}}
\usepackage{bookmark}
\IfFileExists{xurl.sty}{\usepackage{xurl}}{} % add URL line breaks if available
\urlstyle{same}
% fallback for those not using the hyperref driver hyperxmp:
\makeatletter
\@ifundefined{xmpquote}{\newcommand{\xmpquote}[1]{#1}}{}
\makeatother
\hypersetup{
  hidelinks,
  pdfcreator={LaTeX via pandoc}}

\author{}
\date{}

\begin{document}

{
\setcounter{tocdepth}{2}
\tableofcontents
}
\section{WebCoder --- Ghidul de
securitate}\label{webcoder--ghidul-de-securitate}

\begin{quote}
Rezumat într-un paragraf: Postura de securitate a WebCoder stă pe cinci
piloni --- un sandbox de evaluare izolat, atins doar printr-un proxy
Docker de filtrare, autentificare pe cookie-uri httpOnly cu rotire și
CSRF, limite de rate stratificate pe fiecare cale de scriere, un sistem
de moderare/raportare cu audit trail și securitate operațională onestă
(backupuri, exerciții de restore, throttle-uri publicate, chitanțe de
sandbox). Ghidul documentează fiecare control, unde locuiește în cod și
cum îl verifici.
\end{quote}

\textbf{Cui se adresează:} operatori, reviewer-i de securitate,
profesori/părinți care întreabă „cum rulează codul elevului meu?".
Perechea în engleză:
\href{../guide-en/SECURITY.md}{\texttt{docs/guide-en/SECURITY.md}}.

\begin{center}\rule{0.5\linewidth}{0.5pt}\end{center}

\subsection{1. Sandbox-ul evaluatorului (cod
neverificat)}\label{1-sandbox-ul-evaluatorului-cod-neverificat}

Orice execuție judecată rulează în containerul \texttt{judge-runner} cu:

\begin{itemize}
\tightlist
\item
  \texttt{-\/-network\ none} --- fără rețea în timpul compilării sau
  execuției (verificat în CI);
\item
  mount read-only al submissionului (scriibil doar prin calea
  orchestrată \texttt{/srv/judge});
\item
  toate capabilitățile Linux eliminate, \texttt{no-new-privileges}, uid
  1000, limite de memorie/CPU/PID;
\item
  limite de timp wall+CPU și max-RSS per test, persistate ca chitanțe pe
  fiecare rând de rezultat.
\end{itemize}

Worker-ul ajunge la daemon prin \textbf{socket-proxy}
(tecnativa/docker-socket-proxy) care acordă doar containers/exec POST
(+ping/version): \texttt{docker\ run\ hello-world} din backend răspunde
403; \texttt{docker\ exec\ webcoder-judge-runner\ id} reușește. Smoke-ul
de acceptare e în RUNBOOK §socket-proxy.

\textbf{Biblioteci custom} (python3
\texttt{custom\_libraries\_allowed}): pip install-ul rulează într-un
container one-shot pe o rețea internă al cărei unic egress e un proxy
allowlist PyPI pre-populat; liniile de cerință trebuie să se potrivească
cu \texttt{nume} / \texttt{nume==versiune} (git+/URL-uri/căi → 400 la
submit). Postura implicită e fail-closed.

\textbf{Chitanțele de sandbox} --- \texttt{/sandbox-receipts} (orientat
spre utilizator): limitele și rezultatele fiecărei rulări
scratch/stress, auditabile.

\subsection{2. Autentificarea și
sesiunile}\label{2-autentificarea-ux219i-sesiunile}

\begin{itemize}
\tightlist
\item
  JWT access (5 min) + refresh (24 h) stau în \textbf{cookie-uri
  httpOnly} (\texttt{wc\_access}, \texttt{wc\_refresh}, SameSite=Lax în
  dev / None+Secure cross-site); rotirea blacklistează refresh-ul
  anterior; logout-ul blacklistează; nimic de forma token nu atinge
  \texttt{localStorage}.
\item
  Metodele unsafe cer header-ul CSRF double-submit
  (\texttt{GET\ /api/v1/auth/csrf/} eliberează cookie-ul).
\item
  Înregistrarea nu poate auto-atribui roluri (kill-ul B1); schimbările
  de rol sunt doar de admin.
\item
  2FA (TOTP) pentru rolurile privilegiată --- enforce-ment progresiv
  prin \texttt{TWO\_FACTOR\_ENFORCE}.
\item
  Token-uri API personale pentru API-ul public read-only (revocabile;
  tier de throttle mai ridicat).
\item
  Resetarea parolei prin token-uri single-use, conștiente de email;
  providerii OAuth cu scope minim (GitHub \texttt{user,user:email} în
  profilul de produs).
\end{itemize}

\subsection{3. Limitarea ratei ( suprafața de abuz și
DoS)}\label{3-limitarea-ratei--suprafaux21ba-de-abuz-ux219i-dos}

Pe redis, partajată între procese (env-urile \texttt{THROTTLE\_*}):

{\def\LTcaptype{none} % do not increment counter
\begin{longtable}[]{@{}ll@{}}
\toprule\noalign{}
Scope & Implicit \\
\midrule\noalign{}
\endhead
\bottomrule\noalign{}
\endlastfoot
baza anonimă & 60/min \\
baza autentificată & 240/min \\
auth (login/înregistrare) & 10/min \\
creare submission & 20/oră \\
scratch / guest try-it / sample-run & capuri proprii (guest 5/h per
IP) \\
API public & 60/min (240/min cu token) \\
propuneri de traducere & 10/h \\
\end{longtable}
}

\subsection{4. Siguranța conținutului și a
comunității}\label{4-siguranux21ba-conux21binutului-ux219i-a-comunitux103ux21bii}

\begin{itemize}
\tightlist
\item
  Butoane de raportare pe postări, mesaje, probleme; coadă de moderare
  cu tracking de SLA și audit trail; proces de apeluri.
\item
  Mesageria: raportare per mesaj, filtru de conținut RO+EN de tip
  marcare-nu-blocare, blocări end-to-end, capace de trimitere (30/h, 10
  thread-uri/h), politica de contact cu minorii, transcripturi vizibile
  doar la raportare.
\item
  Enforce-ment spoiler server-side în thread-urile de problemă;
  ascundere-la-prag cu comutator; vot ponderat cu încredere.
\item
  Minorii (GDPR Art. 8, România 16+): profiluri ne-publice implicit,
  județ/școală ascunse, înregistrarea consimțământului parental,
  înscriere mediată de școală.
\end{itemize}

\subsection{5. Protecția datelor}\label{5-protecux21bia-datelor}

\begin{itemize}
\tightlist
\item
  GDPR: endpoint de export de date, soft-delete + purgare programată,
  notificare de cookie-uri, paginile \texttt{/terms} +
  \texttt{/privacy}, șablon DPA pentru școli
  (\texttt{docs/DPA-TEMPLATE.md}).
\item
  Outputurile per test: outputurile actual/expected ale testelor ascunse
  se strip-ează pentru privitorii neprivilegiați; datele demo sunt
  publice prin design; task-ul de retenție golește outputurile vechi
  după 90 de zile (verdictele/timpurile se păstrează mereu).
\item
  Backupuri: sidecar \texttt{pg\_dump\ -Fc} zilnic cu retenție 7z/4s/6l;
  copia în afara hostului e decizie de operator; procedura de restore +
  verificare cu accepțiune pe contoarele de rânduri.
\end{itemize}

\subsection{6. Securitate
operațională}\label{6-securitate-operaux21bionalux103}

\begin{itemize}
\tightlist
\item
  Secrete: compose-ul e fail-loud (\texttt{\$\{SECRET\_KEY?err\}});
  profilul de produs (\texttt{webcoder\_api.settings\_prod}) e deny-all
  până se setează
  \texttt{ALLOWED\_HOSTS}/\texttt{CORS\_ALLOWED\_ORIGINS}; HSTS,
  cookie-uri secure, SSL redirect, referrer policy, nosniff.
\item
  Postura de dependențe: Dependabot + CodeQL + configurare OSV scanner
  în repo.
\item
  Observabilitate: \texttt{/healthz} (liveness) + \texttt{/readyz}
  (db/cache/broker), loguri JSON, metrici de admin.
\item
  Contact de incident: \texttt{/security} documentează raportarea;
  pagina de bug-bounty formalizează scope-ul.
\end{itemize}

\subsection{7. Cum verifici
afirmațiile}\label{7-cum-verifici-afirmaux21biile}

\begin{Shaded}
\begin{Highlighting}[]
\CommentTok{\# izolarea de rețea a sandbox{-}ului}
\ExtensionTok{docker}\NormalTok{ exec webcoder{-}judge{-}runner python }\AttributeTok{{-}c} \StringTok{"import urllib.request;urllib.request.urlopen(\textquotesingle{}http://p.org\textquotesingle{})"}   \CommentTok{\# trebuie să eșueze}

\CommentTok{\# socket{-}proxy refuză container{-}create}
\ExtensionTok{docker}\NormalTok{ compose exec backend docker }\AttributeTok{{-}H}\NormalTok{ tcp://socket{-}proxy:2375 run }\AttributeTok{{-}{-}rm}\NormalTok{ hello{-}world   }\CommentTok{\# 403}

\CommentTok{\# un submission judecat rămâne AC (load smoke cu round{-}trip real)}
\ExtensionTok{python}\NormalTok{ scripts/load\_smoke.py }\AttributeTok{{-}{-}base}\NormalTok{ http://127.0.0.1:8010 }\AttributeTok{{-}{-}judge}

\CommentTok{\# exercițiul de restore}
\ExtensionTok{scripts/restore{-}drill.sh}
\end{Highlighting}
\end{Shaded}

Context arhitectural complet:
\href{ARHITECTURA.md}{\texttt{ARHITECTURA.md}}. Runbook operațional:
\texttt{docs/RUNBOOK.md}.

\end{document}
